% !TeX spellcheck = es_ES
\documentclass[letterpaper,11pt]{article}
\usepackage[utf8]{inputenc}
\usepackage[T1]{fontenc}
\usepackage[spanish]{babel}
\decimalpoint % separador decimal con punto, como en las tablas de la fuente original
\usepackage{lmodern,microtype}
\usepackage{amsmath,amssymb,amsfonts}
\usepackage{geometry}
\geometry{margin=2.54cm}
\usepackage{setspace}
\setstretch{1.15}
\usepackage{enumitem}
\usepackage{booktabs}
\usepackage{color}
\usepackage{hyperref}
\definecolor{ntr}{rgb}{0.8,0,0}
\hypersetup{colorlinks=true,linkcolor=ntr,citecolor=ntr,urlcolor=black}

% Portada (Luis Cabral): titulo \LARGE sans-serif alineado a la izquierda, sin numero de pagina
\makeatletter
    \renewcommand{\@maketitle}{%
    \newpage\parindent0in\null\vskip2em%
    {\LARGE\textbf{\textsf{\@title}}\par\vskip2em}%
    \@author\par\vskip2em%
    \@date%
    \par
    \thispagestyle{empty}\setcounter{page}{1}%
    }%
\makeatother

\DeclareMathOperator{\Var}{Var}
\DeclareMathOperator{\Cov}{Cov}
\DeclareMathOperator{\Corr}{Corr}
\DeclareMathOperator{\E}{\mathbb{E}}
\DeclareMathOperator{\N}{Normal}
\DeclareMathOperator{\LP}{\mathbb{L}}
\newcommand\independent{\protect\mathpalette{\protect\independenT}{\perp}}
\def\independenT#1#2{\mathrel{\rlap{$#1#2$}\mkern2mu{#1#2}}}
% Teoria asintotica
\newcommand{\SR}{\mathbb{R}}
\DeclareMathOperator{\plim}{plim}
\DeclareMathOperator{\Avar}{AVar}
\newcommand{\pto}{\stackrel{p}{\longrightarrow}}
\newcommand{\dto}{\stackrel{d}{\longrightarrow}}
\newcommand{\adistr}{\stackrel{a}{\sim}}
\newcommand{\epsi}{\varepsilon}


\title{Análisis Econométrico \\ Ayudantía 5}
\author{Magíster en Economía, Universidad Alberto Hurtado}
\date{}

% Funciones de R sugeridas: entre corchetes, alineadas a la derecha; si no caben en la
% linea pasan enteras a la siguiente.
\newcommand{\fn}[1]{{\unskip\nobreak\hfill\penalty50\hskip1em\null\nobreak\hfill
    \mbox{[#1]}\parfillskip=0pt\finalhyphendemerits=0\par}}

\begin{document}
\maketitle

\noindent Temas: primeros pasos en R. Importar datos, crear variables, estadísticas
descriptivas, gráficos, regresión simple y múltiple por MCO, y tablas de resultados.

\medskip
\noindent En esta ayudantía el ayudante escribe y corre el código en pantalla, ejercicio por
ejercicio; no hace falta traer computador. En cada parte se indican entre corchetes las
funciones de R que se usan. Los ejercicios 1 a 6 son los de la sesión (80 minutos); el 7 queda
como práctica adicional.

\section*{Si quiere replicarlo en su computador}

\noindent Si quiere, puede replicar todo en su computador, durante la ayudantía o después. El
script completo, \texttt{TA5.R}, y los datos están en el repositorio del curso. Esto es lo que
tiene que hacer:

\begin{enumerate}[itemsep=0.4em]
    \item Instale \href{https://cran.r-project.org}{R} y
          \href{https://posit.co/download/rstudio-desktop/}{RStudio}. Si no puede instalarlos,
          cree una cuenta gratuita en \href{https://posit.cloud}{Posit Cloud}, que corre RStudio
          en el navegador: \emph{New Project $\to$ New Project from Git Repository} y pegue
          \texttt{https://github.com/ramirode/econometria\_mae.git}.
    \item Abra RStudio y corra en la consola (una sola vez; tarda varios minutos):
\begin{verbatim}
install.packages(c("tidyverse", "haven", "sandwich",
                   "modelsummary", "kableExtra"))
\end{verbatim}
    \item Cree una carpeta \texttt{TA5} con dos subcarpetas, \texttt{code} e \texttt{input}.
          Guarde el script \texttt{TA5.R} en \texttt{code} y los datos \texttt{ch8\_cps.dta} en
          \texttt{input}. Ambos archivos están en el repositorio del curso, en
          \href{https://github.com/ramirode/econometria_mae/tree/main/TA/TA5}{\texttt{TA/TA5}}.
          En Posit Cloud la carpeta ya existe: \texttt{/cloud/project/TA/TA5}.
    \item Abra \texttt{TA5.R} y cambie la línea \texttt{dir\_proyecto <- "..."} por la ruta a
          su carpeta \texttt{TA5}. Es lo único que hay que modificar.
\end{enumerate}

\section*{Los datos}

\noindent El archivo \texttt{ch8\_cps.dta} contiene datos de trabajadores de tiempo completo de
entre 21 y 64 años encuestados en la \emph{Current Population Survey} de EEUU de marzo de 2005.
Las variables son:
\begin{itemize}[itemsep=0em]
    \item \texttt{ahe}: salario promedio por hora en 2004 (en dólares),
    \item \texttt{yrseduc}: años de educación,
    \item \texttt{female}: igual a 1 si la persona es mujer,
    \item \texttt{age}: edad en años,
    \item \texttt{northeast}, \texttt{midwest}, \texttt{south}, \texttt{west}: dummies de región.
\end{itemize}

\section*{Ejercicios}

\begin{enumerate}[itemsep=0.8em]

%% ---------------------------------------------------------------- 1
\item \emph{Importar e inspeccionar los datos.}
    \begin{enumerate}
        \item Fije la carpeta de trabajo en \texttt{TA5} y cargue los paquetes.
              \fn{\texttt{setwd}, \texttt{library}}
        \item Importe los datos y guárdelos en un objeto llamado \texttt{d}.
              \fn{\texttt{read\_dta}}
        \item ¿Cuántas observaciones y cuántas variables tiene la base? ¿Qué es una observación?
              \fn{\texttt{nrow}, \texttt{ncol}, \texttt{names}, \texttt{head}, \texttt{glimpse}}
        \item Revise el mínimo, la media y el máximo de cada variable. ¿Hay valores faltantes o
              valores que le llamen la atención?
              \fn{\texttt{summary}}
    \end{enumerate}

%% ---------------------------------------------------------------- 2
\item \emph{Crear variables.}
    \begin{enumerate}
        \item Cree la variable \texttt{lahe}, el logaritmo natural de \texttt{ahe}.
              \fn{\texttt{mutate}, \texttt{log}}
        \item Cree la variable \texttt{sexo}, igual a \texttt{"Mujer"} si \texttt{female} es 1 y
              a \texttt{"Hombre"} si es 0. Verifique que quedó bien construida.
              \fn{\texttt{if\_else}, \texttt{count}}
    \end{enumerate}

%% ---------------------------------------------------------------- 3
\item \emph{Estadísticas descriptivas.}
    \begin{enumerate}
        \item Calcule la media y la desviación estándar de \texttt{ahe}, \texttt{lahe},
              \texttt{yrseduc}, \texttt{age} y \texttt{female}. ¿Cómo se interpreta la media de
              \texttt{female}?
              \fn{\texttt{summarise}, \texttt{mean}, \texttt{sd}}
        \item Calcule el número de observaciones y las medias de \texttt{ahe}, \texttt{yrseduc} y
              \texttt{age} por separado para hombres y mujeres. ¿Cuál es la brecha de salarios
              entre hombres y mujeres, en porcentaje? ¿Quiénes tienen más educación?
              \fn{\texttt{group\_by}, \texttt{n}}
        \item Construya una tabla de frecuencias de \texttt{yrseduc}. ¿Cuáles son los niveles de
              educación más frecuentes?
              \fn{\texttt{count}}
    \end{enumerate}

%% ---------------------------------------------------------------- 4
\item \emph{Gráficos.}
    \begin{enumerate}
        \item Grafique un histograma de \texttt{ahe} y otro de \texttt{lahe}. Describa la forma
              de cada distribución.
              \fn{\texttt{ggplot}, \texttt{geom\_histogram}}
        \item Compare la distribución de \texttt{lahe} de hombres y mujeres con un diagrama de
              caja.
              \fn{\texttt{geom\_boxplot}}
        \item Calcule la media de \texttt{lahe} para cada valor de \texttt{yrseduc} y grafíquela
              contra \texttt{yrseduc}. ¿Qué objeto poblacional está estimando? ¿Le parece
              aproximadamente lineal?
              \fn{\texttt{group\_by}, \texttt{geom\_point}}
        \item Guarde los gráficos como imágenes en una carpeta \texttt{output}.
              \fn{\texttt{dir.create}, \texttt{ggsave}}
    \end{enumerate}

%% ---------------------------------------------------------------- 5
\item \emph{Regresión simple.} Considere el modelo
    \begin{align*}
    lahe_i = \beta_0 + \beta_1\, yrseduc_i + u_i.
    \end{align*}
    \begin{enumerate}
        \item Estime el modelo por MCO. Interprete $\hat\beta_1$.
              \fn{\texttt{lm}, \texttt{summary}, \texttt{coef}}
        \item Calcule los errores estándar robustos a heterocedasticidad y compárelos con los
              que reporta \texttt{summary}.
              \fn{\texttt{vcovHC}, \texttt{diag}, \texttt{sqrt}}
        \item Verifique que $\hat\beta_1$ es igual a la covarianza muestral entre
              \texttt{yrseduc} y \texttt{lahe} dividida por la varianza muestral de
              \texttt{yrseduc}, y que $\hat\beta_0=\overline{lahe}-\hat\beta_1\,\overline{yrseduc}$.
              \fn{\texttt{cov}, \texttt{var}}
        \item Guarde los valores predichos y los residuos como variables nuevas. Verifique que
              los residuos tienen media cero y covarianza cero con \texttt{yrseduc}. ¿Por qué
              se cumple esto siempre?
              \fn{\texttt{fitted}, \texttt{resid}}
        \item Agregue la recta de regresión al gráfico del ejercicio 4(c). ¿Qué log salario
              predice el modelo para una persona con 12 años de educación? ¿Y con 16?
              \fn{\texttt{geom\_smooth}, \texttt{predict}}
    \end{enumerate}

%% ---------------------------------------------------------------- 6
\item \emph{Regresión múltiple y tabla de resultados.} Considere los modelos
    \begin{align*}
    lahe_i &= \beta_0 + \beta_1\, yrseduc_i + \beta_2\, female_i + u_i,\\
    lahe_i &= \beta_0 + \beta_1\, yrseduc_i + \beta_2\, female_i + \beta_3\, age_i + u_i.
    \end{align*}
    \begin{enumerate}
        \item Estime los dos modelos. Interprete $\hat\beta_2$ y compárelo con la brecha de
              salarios del ejercicio 3(b).
              \fn{\texttt{lm}}
        \item El coeficiente de \texttt{yrseduc}, ¿aumenta o disminuye al agregar
              \texttt{female}? Explique el signo del cambio con la fórmula del sesgo de variable
              omitida y verifique que la fórmula se cumple exactamente en la muestra.
        \item Arme una tabla con los tres modelos (el del ejercicio 5 y estos dos), con errores
              estándar robustos, y guárdela en \texttt{output}.
              \fn{\texttt{list}, \texttt{modelsummary}}
        \item ¿Cree que $\hat\beta_1$ mide el efecto causal de un año más de educación sobre el
              salario? Mencione una variable omitida y el signo del sesgo que genera.
    \end{enumerate}

\end{enumerate}

\section*{Práctica adicional}

\begin{enumerate}[itemsep=1.2em,start=7]
\item \emph{MCO con matrices y regresiones por sexo.} Sea $\mathbf{X}$ la matriz de $N\times 4$ cuyas columnas son un
    vector de unos, \texttt{yrseduc}, \texttt{female} y \texttt{age}, y sea $\mathbf{y}$ el
    vector de \texttt{lahe}.
    \begin{enumerate}
        \item Calcule $\hat\beta=(\mathbf{X}'\mathbf{X})^{-1}\mathbf{X}'\mathbf{y}$ y verifique
              que coincide con el resultado de \texttt{lm}.
              \fn{\texttt{cbind}, \texttt{t}, \texttt{\%*\%}, \texttt{solve}}
        \item Estime la regresión simple del ejercicio 5 por separado para hombres y para
              mujeres. ¿Para quiénes es mayor el retorno a la educación?
              \fn{\texttt{filter}}
    \end{enumerate}
\end{enumerate}

\end{document}
